\documentclass[11pt]{article}

\usepackage[T1]{fontenc}
\usepackage[utf8]{inputenc}
\usepackage{lmodern}
\usepackage[english]{babel}
\usepackage{microtype}
\usepackage{geometry}
\geometry{margin=1in}
\usepackage{amsmath,amssymb,amsthm,mathtools}
\usepackage{booktabs,array,longtable}
\usepackage{enumitem}
\usepackage{xcolor}
\usepackage{hyperref}
\usepackage{url}
\usepackage{fancyhdr}
\usepackage{setspace}

\hypersetup{
  colorlinks=true,
  linkcolor=blue!55!black,
  citecolor=green!45!black,
  urlcolor=blue!65!black,
  pdftitle={Alaoglu--Erdos: Rational-Orbit Rigidity and a Proof Audit},
  pdfauthor={Research continuation prepared for Vladimir Reshetnikov}
}

\setlength{\parindent}{0pt}
\setlength{\parskip}{0.65em}
\setstretch{1.04}
\pagestyle{fancy}
\setlength{\headheight}{14pt}
\fancyhf{}
\lhead{Alaoglu--Erd\H{o}s continuation}
\rhead{August 21, 2026}
\cfoot{\thepage}

\newtheorem{theorem}{Theorem}[section]
\newtheorem{proposition}[theorem]{Proposition}
\newtheorem{lemma}[theorem]{Lemma}
\newtheorem{corollary}[theorem]{Corollary}
\newtheorem{criterion}[theorem]{Criterion}
\newtheorem{observation}[theorem]{Observation}
\theoremstyle{definition}
\newtheorem{definition}[theorem]{Definition}
\newtheorem{problem}[theorem]{Research problem}
\theoremstyle{remark}
\newtheorem{remark}[theorem]{Remark}
\newtheorem{warning}[theorem]{Audit warning}

\newcommand{\Q}{\mathbb Q}
\newcommand{\R}{\mathbb R}
\newcommand{\Z}{\mathbb Z}
\newcommand{\N}{\mathbb N}
\newcommand{\Nzero}{\mathbb N_0}
\newcommand{\Abar}{\overline{\mathbb Q}}
\newcommand{\fract}[1]{\{#1\}}
\newcommand{\ord}{\operatorname{ord}}
\newcommand{\den}{\operatorname{den}}
\newcommand{\Hgt}{\operatorname{H}}
\newcommand{\dist}{\operatorname{dist}}
\newcommand{\Span}{\operatorname{span}}

\newcommand{\statusbox}[1]{%
  \begin{center}
  \fcolorbox{black}{gray!8}{\parbox{0.92\textwidth}{\textbf{Proof-status statement.} #1}}
  \end{center}
}

\title{\textbf{Alaoglu--Erd\H{o}s After the Unified Report}\\[0.4em]
\Large Rational-Orbit Denominators, Farey Entropy, Smooth Rigidity,\\
and an Audited Attempt at the Remaining Proof}
\author{Research continuation prepared for Vladimir Reshetnikov}
\date{August 21, 2026}

\begin{document}
\maketitle

\begin{abstract}
The Alaoglu--Erd\H{o}s conjecture asserts that a real number $x$ satisfying
$2^x,3^x\in\Z$ must be an integer.  The supplied unified report is already a very large
research dossier: it develops exact conditional structure, determinant methods, matrix and
operator reformulations, $p$-adic diagnostics, approximation barriers, and extensive
formal verification.  The present document is a separate continuation rather than a
restatement of that dossier.

The new work concentrates on a route suggested by the July 2026 theorem of Lelis,
Moreira, and Silva on smooth functions mapping rationals to rationals with controlled
denominators.  Assuming a counterexample, we first sharpen the compact-orbit description to
an exact numerator, denominator, and height law.  We then prove a new
\emph{Farey-to-orbit entropy theorem}: every injective rational parametrization of the dense
conditional orbit must distort denominators at least like $\exp(cQ^2)$ on rationals of
height at most $Q$.  Consequently, no polynomially tame reparametrization can place the
orbit in the regime of known rational-mapping rigidity theorems.  We also derive an
orbit-adapted subunit-locking inequality and show quantitatively why every fixed-order
rational-contact or divided-difference certificate is underpowered: its analytic gain is
only polynomial in the orbit index, whereas its arithmetic denominator cost is exponential.

These results rigorously close several plausible proof routes and identify a precise missing
amplification principle.  They do not supply an unconditional proof of the conjecture.
Current transcendence theory still places the decisive $2\times2$ logarithmic determinant at
the Four Exponentials boundary; the recent Matrix Coefficient Conjecture is equivalent to
Four Exponentials in size two and remains conjectural.  No statement in this report is
promoted beyond the proof actually given.
\end{abstract}

\statusbox{The report proves all numbered elementary and height-theoretic results below,
including the exact orbit-height formula and the Farey-to-orbit entropy theorem.  It does
\emph{not} relabel an unresolved implication as a proof of the conjecture.  The final
size-two logarithmic step remains open in the verified literature; every conditional input
is marked explicitly.}

\tableofcontents
\newpage

\section{Scope, non-duplication, and the logical standard}

\subsection{What was taken from the supplied report}

The source file
\begin{center}
\texttt{/Article/alaoglu-erdos-unified-report/alaoglu-erdos-unified-report.tex}
\end{center}
contains roughly nineteen thousand lines and already consolidates a large collection of
independent continuations.  In particular, it contains:

\begin{itemize}[leftmargin=2em]
\item the rational-exponent reduction and the Gelfond--Schneider consequence;
\item the exact logarithmic-determinant and Four Exponentials formulation;
\item an exact conditional monoid of solutions, a primitive generator, and compact
      synchronized dyadic--triadic orbits;
\item determinant, generalized Vandermonde, HCIZ, Smith-divisor, finite-difference,
      approximation, matrix, spectral, Mahler, and operator analyses;
\item exact reduced denominator exponents, qualitative and asymptotic compact-height
      counting, and rational-contact transversality;
\item an extensive catalogue of failure points and a literature audit through most of 2026.
\end{itemize}

Those arguments are not reproduced.  We use only a minimal baseline and independently prove
the refinements needed here.  The genuinely new core of this continuation is:

\begin{enumerate}[leftmargin=2.2em]
\item an exact closed formula for the heights of every point in the compact conditional
      orbit, not merely an $\asymp$ estimate;
\item a Farey-counting theorem showing that any injection from all rationals in an interval
      into that orbit has denominator distortion $\exp(cQ^2)$;
\item a corresponding no-go theorem for polynomially tame rational reparametrizations and
      for direct use of the 2026 smooth rational-map rigidity theorem;
\item an orbit-adapted comparison between arithmetic lower bounds and analytic contact
      estimates, isolating the exact exponential-versus-polynomial scale mismatch;
\item a set of proof-producing targets stated strongly enough that crossing any one of them
      would be a genuine advance rather than another reformulation.
\end{enumerate}

\subsection{Current status of the transcendence input}

Dasgupta and Kakde formulate the Matrix Coefficient Conjecture for matrices of logarithms of
algebraic numbers and explicitly note that in dimension two it is equivalent to the Four
Exponentials Conjecture \cite{DasguptaKakde2024}.  Their paper proves several implications and
builds a support-sensitive strategy, but leaves a key implication open.  Thus it cannot be
cited as a proof of the determinant nonvanishing needed here.

The recent paper of Lelis, Moreira, and Silva proves a powerful rigidity theorem for
$C^{2k+1}$ functions mapping every rational in a neighborhood to a rational with denominator
$O(q^k)$, provided a rational jet through order $2k$ is available
\cite{LelisMoreiraSilva2026}.  This theorem is highly relevant to the compact rational orbit,
but two hypotheses fail in a sharp way: the conditional graph contains only logarithmically
many points of height at most $H$, and its first derivative at every rational graph point is
transcendental.  Sections~\ref{sec:farey}--\ref{sec:lelis} quantify the first failure and make
the second explicit.

\subsection{What counts as a proof}

A proof must derive a contradiction from only
\[
  x\in\R,\qquad 2^x\in\Z,\qquad 3^x\in\Z,
\]
using established theorems.  The following are not proofs:

\begin{itemize}[leftmargin=2em]
\item assuming Four Exponentials, Schanuel, or the Matrix Coefficient Conjecture without
      labeling the assumption;
\item replacing exact rationality by numerical near-rationality;
\item invoking density of a rational orbit as though a continuous function mapping a dense
      rational set to rationals had to be rational;
\item using an arithmetic lower bound $e^{-cN}$ against an analytic upper bound $N^{-d}$ and
      reversing their asymptotic order;
\item treating real logarithmic proportionality as a $p$-adic proportionality without first
      proving continuity of the induced homomorphism.
\end{itemize}

The scale comparison in this report was designed specifically to expose these errors before
they enter a purported proof.

\section{Minimal reduction to the exact logarithmic wall}
\label{sec:baseline}

\begin{proposition}[Elementary reduction]\label{prop:elementary}
Suppose $x\in\R$ and $2^x,3^x\in\Z$.  Then $x\ge0$.  If $x\in\Q$, then $x\in\Nzero$.
If $x\notin\Z$, then $x$ is transcendental.
\end{proposition}

\begin{proof}
If $x<0$, then $0<2^x<1$, so $2^x$ is not an integer.  Now write
$x=a/b$ in lowest terms with $b>0$.  If $2^{a/b}=m\in\Z$, then
$2^a=m^b$.  Taking the $2$-adic valuation gives
$a=bv_2(m)$, so $b\mid a$.  Coprimality forces $b=1$.

If $x$ were algebraic and irrational, Gelfond--Schneider would make $2^x$
transcendental, contrary to $2^x\in\Z$.  Hence every nonintegral candidate is
transcendental.
\end{proof}

Put
\[
  m=2^x,\qquad n=3^x.
\]
Then
\begin{equation}
  \log m\,\log 3-\log n\,\log 2=0.
  \label{eq:logdet}
\end{equation}
Equivalently,
\[
 \det\begin{pmatrix}
 \log2&\log3\\
 \log m&\log n
 \end{pmatrix}=0.
\]
Since $\log m=x\log2$ and $\log n=x\log3$, this determinant vanishes tautologically under
the hypotheses.  The issue is whether a singular matrix of logarithms of algebraic numbers
can have both row and column directions rationally independent.

\begin{proposition}[Conditional Four Exponentials proof]\label{prop:four-exp}
The Four Exponentials Conjecture implies the Alaoglu--Erd\H{o}s conjecture.
\end{proposition}

\begin{proof}
Assume $x\notin\Z$.  Proposition~\ref{prop:elementary} gives $x\notin\Q$, so
$1,x$ are linearly independent over $\Q$.  The numbers $\log2,\log3$ are also linearly
independent over $\Q$, since $2^a3^b=1$ with $a,b\in\Z$ implies $a=b=0$.  The four
exponentials
\[
 e^{\log2}=2,\quad e^{\log3}=3,\quad e^{x\log2}=m,\quad e^{x\log3}=n
\]
are all algebraic.  Four Exponentials predicts that at least one must be transcendental, a
contradiction.
\end{proof}

This proposition is not an unconditional proof because Four Exponentials is open.  The rest
of the report asks whether the stronger fact that all four exponentials are positive
integers supplies an independent route around that conjecture.

\section{The compact conditional orbit with exact heights}
\label{sec:orbit}

Assume for the remainder of Sections~\ref{sec:orbit}--\ref{sec:contact} that a nonintegral
counterexample exists.  Write
\[
  x=B+\alpha,
  \qquad B=\lfloor x\rfloor\in\Nzero,
  \qquad 0<\alpha<1.
\]
By Proposition~\ref{prop:elementary}, $\alpha$ is irrational.  Set
\[
  U=2^\alpha=\frac{2^x}{2^B}\in\Q\cap(1,2),
  \qquad
  V=3^\alpha=\frac{3^x}{3^B}\in\Q\cap(1,3).
\]
Write these rationals in lowest terms as
\begin{equation}
  U=\frac{W}{2^r},\qquad W>1\text{ odd},\quad r\ge1,
  \label{eq:Ured}
\end{equation}
\begin{equation}
  V=\frac{Z}{3^s},\qquad Z>1,\quad 3\nmid Z,\quad s\ge0.
  \label{eq:Vred}
\end{equation}
The denominator restrictions follow because $U$ is an integer divided by a power of $2$,
and similarly for $V$.  Moreover,
\begin{equation}
  r+\alpha=\log_2 W,
  \qquad
  s+\alpha=\log_3 Z.
  \label{eq:slope-two-forms}
\end{equation}

For $k\ge0$, define
\begin{equation}
  R_k=2^{\fract{k\alpha}},
  \qquad
  S_k=3^{\fract{k\alpha}}.
  \label{eq:compact-orbit}
\end{equation}
The set $\{\fract{k\alpha}:k\ge0\}$ is dense in $[0,1)$, so
$(R_k,S_k)$ is dense on the analytic arc
\[
  \mathcal A=\{(2^t,3^t):0\le t<1\}.
\]
The next theorem gives exact arithmetic data for every orbit point.

\begin{theorem}[Exact reduced orbit law]\label{thm:exact-orbit}
For every integer $k\ge1$, the following are lowest-term expressions:
\begin{align}
 R_k&=\frac{W^k}{2^{\lfloor k\log_2W\rfloor}},
 \label{eq:Rexact}\\
 S_k&=\frac{Z^k}{3^{\lfloor k\log_3Z\rfloor}}.
 \label{eq:Sexact}
\end{align}
Consequently,
\begin{align}
 \Hgt(R_k)&=W^k, & \den(R_k)&=2^{\lfloor k\log_2W\rfloor},
 \label{eq:Rheight}\\
 \Hgt(S_k)&=Z^k, & \den(S_k)&=3^{\lfloor k\log_3Z\rfloor}.
 \label{eq:Sheight}
\end{align}
In particular,
\begin{align}
 \frac{W^k}{2}<\den(R_k)<W^k,
 \qquad
 \frac{Z^k}{3}<\den(S_k)<Z^k.
 \label{eq:den-comparable}
\end{align}
\end{theorem}

\begin{proof}
From $U=W/2^r=2^\alpha$,
\[
 R_k=2^{k\alpha-\lfloor k\alpha\rfloor}
     =\frac{W^k}{2^{rk+\lfloor k\alpha\rfloor}}.
\]
Equation~\eqref{eq:slope-two-forms} gives
\[
 rk+\lfloor k\alpha\rfloor
 =\lfloor k(r+\alpha)\rfloor
 =\lfloor k\log_2W\rfloor.
\]
The numerator is odd, so the fraction is reduced.  The proof of
\eqref{eq:Sexact} is identical, using $3\nmid Z$.

Because $1<R_k<2$, the numerator of the reduced fraction is larger than its denominator and
smaller than twice its denominator.  Hence its projective height is exactly $W^k$, and
$W^k/2<\den(R_k)<W^k$.  The base-$3$ statement follows from $1<S_k<3$.
\end{proof}

\begin{corollary}[Exact compact-orbit counting]\label{cor:exact-count}
Let
\[
  H_k=\max\{\Hgt(R_k),\Hgt(S_k)\},
  \qquad C_0=\max\{W,Z\}.
\]
Then $H_k=C_0^k$.  Hence, for $H\ge1$,
\begin{equation}
 \#\{k\ge0:H_k\le H\}
 =1+\left\lfloor\frac{\log H}{\log C_0}\right\rfloor.
 \label{eq:exact-count}
\end{equation}
\end{corollary}

\begin{proof}
The height identity follows from Theorem~\ref{thm:exact-orbit}.  The points are distinct
because $\alpha$ is irrational.  The inequality $C_0^k\le H$ is equivalent to
$k\le\log H/\log C_0$.
\end{proof}

The supplied report proves the qualitative dichotomy $1$ versus $\asymp\log H$ for the
compact arc.  Corollary~\ref{cor:exact-count} records the precise conditional counting
function and is the input needed for the entropy argument below.

\section{The Farey-to-orbit entropy mismatch}
\label{sec:farey}

The natural hope inspired by rational-map rigidity is to reparametrize the dense orbit by all
rationals in an interval, then apply a theorem about smooth functions preserving rationality.
This section proves that every such reparametrization is arithmetically violent.

Fix a nondegenerate closed interval $I=[a,b]\subset(0,1)$ and write $|I|=b-a$.  Let
\[
 \mathcal F_I(Q)=\left\{\frac pq\in I:
 1\le q\le Q,\ \gcd(p,q)=1\right\}.
\]
The classical Farey estimate is
\begin{equation}
 \#\mathcal F_I(Q)=\frac{3|I|}{\pi^2}Q^2+O_I(Q\log Q).
 \label{eq:farey}
\end{equation}
A proof sketch is included in Appendix~\ref{app:farey}.

\begin{definition}[Orbit index]
For a point $P=(R_k,S_k)$ in the one-sided compact orbit, write $\nu(P)=k$.
This is well-defined because the orbit points are distinct.
\end{definition}

\begin{theorem}[Farey-to-orbit entropy theorem]\label{thm:farey-orbit}
Let
\[
 \Phi:\Q\cap I\longrightarrow\{(R_k,S_k):k\ge0\}
\]
be injective.  Define
\[
 D_2(Q)=\max_{t\in\mathcal F_I(Q)}\den\bigl(\pi_1\Phi(t)\bigr),
 \qquad
 D_3(Q)=\max_{t\in\mathcal F_I(Q)}\den\bigl(\pi_2\Phi(t)\bigr).
\]
Then
\begin{align}
 \liminf_{Q\to\infty}\frac{\log D_2(Q)}{Q^2}
 &\ge \frac{3|I|}{\pi^2}\log W,
 \label{eq:D2entropy}\\
 \liminf_{Q\to\infty}\frac{\log D_3(Q)}{Q^2}
 &\ge \frac{3|I|}{\pi^2}\log Z.
 \label{eq:D3entropy}
\end{align}
In particular, for every fixed $A>0$, neither $D_2(Q)$ nor $D_3(Q)$ is $O(Q^A)$.
\end{theorem}

\begin{proof}
Put $N_Q=\#\mathcal F_I(Q)$.  The images of these $N_Q$ rationals have distinct orbit
indices.  Therefore at least one image has index at least $N_Q-1$; equivalently,
\[
 \max_{t\in\mathcal F_I(Q)}\nu(\Phi(t))\ge N_Q-1.
\]
By Theorem~\ref{thm:exact-orbit}, an orbit point of index $k$ has first-coordinate
denominator greater than $W^k/2$ and second-coordinate denominator greater than $Z^k/3$.
Consequently,
\[
 D_2(Q)>\frac12W^{N_Q-1},
 \qquad
 D_3(Q)>\frac13Z^{N_Q-1}.
\]
Taking logarithms, dividing by $Q^2$, and using~\eqref{eq:farey} proves the two limits.
Since $\log Q^A/Q^2\to0$, no polynomial bound is possible.
\end{proof}

\begin{remark}
The theorem uses only injectivity and the exact orbit-height law.  It does not use
continuity, differentiability, order preservation, or equidistribution.  Thus adding
analytic regularity to the parametrization cannot remove the arithmetic distortion.
\end{remark}

\begin{corollary}[No polynomially tame rational conjugacy]\label{cor:no-tame-conjugacy}
Let $J=(a,b)$ be a nonempty open interval, and suppose an increasing homeomorphism
$\phi:J\to(1,2)$ satisfies
\[
 \phi(\Q\cap J)=\{R_k:k\ge1\}.
\]
For each rational $t\in J$, let $G(t)=S_k$ where $\phi(t)=R_k$.  Then $G$ is the
restriction to $\Q\cap J$ of the continuous function
\[
 t\longmapsto \phi(t)^{\log_2 3}.
\]
Moreover, for every nondegenerate closed interval $I\Subset J$,
\begin{equation}
 \max_{p/q\in\mathcal F_I(Q)}\den G(p/q)
 \ge \exp\!\left(\left(\frac{3|I|\log Z}{\pi^2}+o(1)\right)Q^2\right).
 \label{eq:conjugacy-distortion}
\end{equation}
Thus no such rational conjugacy can satisfy a denominator estimate
$\den G(p/q)\ll q^A$ for any fixed $A$.
\end{corollary}

\begin{proof}
For a fixed $I\Subset J$, the map $t\mapsto(\phi(t),G(t))$ injects $\Q\cap I$ into the
orbit.  The formula for the continuous extension follows pointwise from
$S_k=R_k^{\log_2 3}$.  Apply Theorem~\ref{thm:farey-orbit} to the second coordinate.
\end{proof}

\begin{corollary}[Bounded multiplicity version]\label{cor:bounded-multiplicity}
If every orbit point has at most $M$ preimages under a map
$\Phi:\Q\cap I\to\{(R_k,S_k)\}$, where $M$ is fixed, then the constants on the right of
\eqref{eq:D2entropy} and~\eqref{eq:D3entropy} may be divided by $M$.
In particular, polynomial denominator growth is still impossible.
\end{corollary}

\begin{proof}
The $N_Q$ input rationals occupy at least $N_Q/M$ distinct orbit indices.  Repeat the proof of
Theorem~\ref{thm:farey-orbit}.
\end{proof}


\begin{proposition}[Prime-level denominator explosion]\label{prop:prime-level}
Let $\Phi:\Q\cap I\to\{(R_k,S_k):k\ge0\}$ be injective.  For every sufficiently large
prime $q$, put
\[
 E_2(q)=\max_{\substack{p/q\in I\\1\le p<q}}\den\bigl(\pi_1\Phi(p/q)\bigr),
 \qquad
 E_3(q)=\max_{\substack{p/q\in I\\1\le p<q}}\den\bigl(\pi_2\Phi(p/q)\bigr).
\]
Then
\begin{align}
 \log E_2(q)&\ge (|I|\log W+o(1))q,\label{eq:prime-E2}\\
 \log E_3(q)&\ge (|I|\log Z+o(1))q.\label{eq:prime-E3}
\end{align}
Thus even when one fixes a single prime denominator $q$, at least one output denominator is
exponential in $q$.
\end{proposition}

\begin{proof}
For prime $q$, every $1\le p<q$ is coprime to $q$, and the number of $p$ with $p/q\in I$ is
$|I|q+O(1)$.  Their images have distinct orbit indices, so one index is at least
$|I|q+O(1)$.  Apply the denominator lower bounds in
Theorem~\ref{thm:exact-orbit}.
\end{proof}

\begin{corollary}[Enumeration form of the entropy mismatch]\label{cor:enumeration}
Enumerate the rationals in $I$ in nondecreasing height and inject them into the orbit.  Among
the first $N$ inputs, some image has each coordinate denominator at least
\[
 \frac12W^{N-1}\quad\text{and}\quad\frac13Z^{N-1},
\]
respectively.  Since the $N$th rational in a height ordering has height $\asymp_I\sqrt N$,
the optimal possible denominator distortion is already $\exp(cH^2)$.
\end{corollary}

\begin{proof}
Distinct images use $N$ distinct nonnegative indices, one of which is at least $N-1$.  The
height assertion follows by inverting the Farey estimate~\eqref{eq:farey}.
\end{proof}

\subsection{Interpretation}

There are $\asymp Q^2$ rationals of height at most $Q$ in any fixed interval, but only
$\asymp\log H$ points of the conditional orbit of height at most $H$.  An injective map from
the first set to the second must therefore solve
\[
 Q^2\lesssim\log H,
\]
forcing $H\gtrsim e^{cQ^2}$.  This is an arithmetic entropy mismatch: the Farey set has
quadratic height entropy, while the cyclic conditional orbit has logarithmic height entropy.

The result is not a contradiction to the existence of the orbit.  It proves instead that
any attempt to ``fill in'' the orbit by a rational parameter must pay a denominator cost far
outside the range where known smooth-rigidity theorems operate.

\section{The 2026 rational-map rigidity theorem and its exact boundary}
\label{sec:lelis}

\subsection{The external theorem}

The following is the form of the recent theorem needed for comparison.

\begin{theorem}[Lelis--Moreira--Silva, 2026; simplified statement]
\label{thm:LMS}
Let $k\ge1$, let $\Omega$ be a connected neighborhood of $0$, and let
$f\in C^{2k+1}(\Omega)$.  Assume
\begin{enumerate}[label=(\roman*),leftmargin=2.2em]
\item $f^{(j)}(0)\in\Q$ for $0\le j\le2k$;
\item $f(\Omega\cap\Q)\subseteq\Q$;
\item $\den f(p/q)\ll q^k$ for rationals $p/q\in\Omega$ with sufficiently large $q$.
\end{enumerate}
Then $f$ is a rational function with rational coefficients of degree at most $k$.
\end{theorem}

The proof combines a rational Pad\'e-type construction from the rational jet with a
subunit-locking argument: after subtraction of the rational approximant, the residual is
$O(q^{-(2k+1)})$ while its denominator is $O(q^{2k})$, so a nonzero rational residual is
impossible for large $q$ \cite{LelisMoreiraSilva2026}.

\subsection{Why the conditional power curve misses both decisive hypotheses}

Put
\[
 \vartheta=\log_2 3=\frac{\log3}{\log2},
 \qquad f(t)=t^\vartheta.
\]
The number $\vartheta$ is transcendental: it is irrational, and if it were algebraic then
Gelfond--Schneider would make $2^\vartheta=3$ transcendental.

Under a counterexample, $f(R_k)=S_k\in\Q$ on a dense set of rational arguments.  This looks
close to Theorem~\ref{thm:LMS}, but it is not close in height geometry.

First, the theorem requires every rational argument in a neighborhood, not a dense cyclic
subset.  The orbit contributes only $O(\log H)$ rational graph points of height at most $H$.
Second, any reparametrization that maps all rationals injectively into the orbit has
$\exp(cq^2)$ denominator distortion by Theorem~\ref{thm:farey-orbit}, rather than the required
$q^k$.

Third, the rational-jet condition fails already in order one.

\begin{proposition}[Transcendental tangent at every rational graph point]
\label{prop:transverse}
Let $a,b\in\Q_{>0}$ satisfy $b=a^\vartheta$.  Then
\[
 f'(a)=\vartheta\frac ba
\]
is transcendental.  Consequently, if $R\in\Q(t)$ is regular at $a$ and $R(a)=b$, then
$f-R$ has a simple zero at $a$; in particular, it cannot have contact order at least two.
\end{proposition}

\begin{proof}
The factor $b/a$ is a nonzero rational and $\vartheta$ is transcendental, so
$\vartheta b/a$ is transcendental.  On the other hand, $R'(a)\in\Q$.  Therefore
$f'(a)-R'(a)\ne0$.
\end{proof}

\begin{corollary}[No direct Lelis--Moreira--Silva reduction]
\label{cor:no-LMS}
Neither of the following can turn a hypothetical Alaoglu--Erd\H{o}s orbit into an
application of Theorem~\ref{thm:LMS}:
\begin{enumerate}[label=(\alph*),leftmargin=2.2em]
\item translating the power function around one of its rational graph points;
\item injectively reparametrizing the entire rational set of a neighborhood by orbit points
      with polynomial denominator growth.
\end{enumerate}
\end{corollary}

\begin{proof}
In (a), the translated function has a transcendental first derivative at the origin by
Proposition~\ref{prop:transverse}.  In (b), polynomial denominator growth contradicts
Theorem~\ref{thm:farey-orbit}.
\end{proof}

\begin{remark}[Density is not the missing theorem]
A continuous function agreeing with $t^\vartheta$ on the dense orbit agrees with it on the
whole interval.  Thus a continuous extension that also mapped every rational to a rational
would force $t^\vartheta\in\Q$ for every rational $t$ in an interval.  That is already
incompatible with the Six Exponentials Theorem after choosing three multiplicatively
independent rational bases.  The hard issue is not uniqueness of the continuous extension;
it is producing rational values away from the rank-two orbit.
\end{remark}

\section{Orbit-adapted subunit locking and the exponential denominator tax}
\label{sec:contact}

The Lelis--Moreira--Silva proof wins because a residual of order $q^{-(2k+1)}$ is a rational
with denominator only $O(q^{2k})$.  On the conditional orbit, the geometric recurrence scale
is polynomial in the orbit index, but the rational denominators are exponential in that
index.  We now make this comparison exact.

\subsection{Return times near the rational point $(1,1)$}

Let $p_j/q_j$ be the continued-fraction convergents of $\alpha$.  Infinitely many are below
$\alpha$, so for an infinite subsequence
\begin{equation}
 0<\delta_j:=q_j\alpha-p_j<\frac1{q_{j+1}}<\frac1{q_j}.
 \label{eq:convergent-return}
\end{equation}
Because $W$ and $2$ are multiplicatively independent, a standard two-logarithm consequence
of Baker's theorem~\cite{Baker1975} gives constants $c_W>0$ and $\kappa_W>1$ such that
\begin{equation}
 |q\alpha-p|\ge c_W q^{-\kappa_W}
 \qquad(q\ge2,\ p\in\Z).
 \label{eq:baker-return-lower}
\end{equation}
Indeed, if $|q\alpha-p|\ge1$ the claim is trivial after shrinking $c_W$; otherwise
$p=O_W(q)$.  After multiplying by $\log2$, the left side is then the nonzero linear
form $q\log W-(rq+p)\log2$, whose integer coefficients have size $O_W(q)$.  In particular,
the return scale in~\eqref{eq:convergent-return} is polynomial, never exponential, in the
orbit index; also $q_{j+1}\ll_W q_j^{\kappa_W}$.

For this subsequence,
\begin{align}
 0<R_{q_j}-1&=2^{\delta_j}-1\ll\frac1{q_j},
 \label{eq:Rreturn}\\
 0<S_{q_j}-1&=3^{\delta_j}-1\ll\frac1{q_j}.
 \label{eq:Sreturn}
\end{align}
The implied constants in these two upper bounds are absolute.

\begin{lemma}[Height of a fixed rational expression]
\label{lem:rational-height}
Let $F\in\Q(X,Y)$ be regular on a neighborhood of the compact arc $\mathcal A$.  There is a
constant $C_F>0$ such that, whenever $F(R_k,S_k)\ne0$,
\begin{equation}
 |F(R_k,S_k)|\ge e^{-C_F k}.
 \label{eq:arith-lower}
\end{equation}
\end{lemma}

\begin{proof}
Write $F=A/B$ with $A,B\in\Z[X,Y]$ after clearing a fixed common denominator.  Let $d$ be a
bound for the total degrees and $H_F$ a bound for the coefficients.  Standard height
inequalities give
\[
 \Hgt(F(R_k,S_k))\le C'_F\Hgt(R_k)^d\Hgt(S_k)^d
 \le C'_F(WZ)^{dk}.
\]
A nonzero rational number of height at most $H$ has absolute value at least $H^{-1}$.
Absorbing the fixed factor $C'_F$ into the exponent proves~\eqref{eq:arith-lower}.
\end{proof}

\begin{theorem}[Fixed algebraic contact has only polynomial strength]
\label{thm:fixed-contact}
Let $0\ne P\in\Q[X,Y]$ satisfy $P(1,1)=0$, and define the finite positive integer
\[
 \sigma(P)=\ord_{t=0}P(2^t,3^t).
\]
Along the subsequence~\eqref{eq:convergent-return}, the values
$P(R_{q_j},S_{q_j})$ are nonzero for all sufficiently large $j$, and there are constants
$c_P,C_P>0$ such that
\begin{equation}
 c_Pq_j^{-\kappa_W\sigma(P)}
 \le |P(R_{q_j},S_{q_j})|
 \le C_Pq_j^{-\sigma(P)}.
 \label{eq:two-sided-contact}
\end{equation}
They also satisfy the weaker height-theoretic estimate
\begin{equation}
 |P(R_{q_j},S_{q_j})|\ge e^{-C'_Pq_j}.
 \label{eq:exp-lower}
\end{equation}
If $P$ has algebraic multiplicity $s$ at $(1,1)$, then $\sigma(P)\ge s$.
Consequently, every fixed polynomial contact certificate has only polynomial smallness in
the orbit index and cannot activate an exponential-denominator subunit argument.
\end{theorem}

\begin{proof}
Put $F(t)=P(2^t,3^t)$.  This analytic function is not identically zero.  Indeed, writing
$P(X,Y)=\sum c_{ab}X^aY^b$ gives
\[
 F(t)=\sum c_{ab}e^{(a\log2+b\log3)t};
\]
the exponents are distinct because $\log2$ and $\log3$ are linearly independent over
$\Q$, and distinct exponential functions are linearly independent.  Hence
$\sigma(P)<\infty$, and for sufficiently small positive $t$,
\[
 c'_P t^{\sigma(P)}\le |F(t)|\le C'_P t^{\sigma(P)}.
\]
Now $P(R_{q_j},S_{q_j})=F(\delta_j)$.  Combining
\eqref{eq:convergent-return} and~\eqref{eq:baker-return-lower} proves
\eqref{eq:two-sided-contact}, and also shows eventual nonvanishing.  The estimate
\eqref{eq:exp-lower} follows independently from Lemma~\ref{lem:rational-height}.
Finally, a Taylor expansion of $P$ in $(X-1,Y-1)$ shows that algebraic multiplicity $s$
forces $F(t)=O(t^s)$, whence $\sigma(P)\ge s$.
\end{proof}

The important point is now two-sided: Baker's theorem prevents exponentially close returns,
and analyticity converts every fixed contact order into a fixed power of that return.  Thus
the actual certificate is polynomially, not exponentially, small.  The elementary rational
lower bound~\eqref{eq:exp-lower} is far too weak to collide with it.

\subsection{A general jet-budget inequality}

The scale comparison can be abstracted.

\begin{proposition}[Orbit jet budget]\label{prop:jet-budget}
Suppose a family of nonzero rational certificates $\Xi_N$ satisfies
\begin{equation}
 \den(\Xi_N)\le \exp(A_NN)
 \label{eq:cert-den}
\end{equation}
and an analytic construction gives
\begin{equation}
 |\Xi_N|\le C_NN^{-d_N}.
 \label{eq:cert-upper}
\end{equation}
A contradiction from the elementary lower bound
$|\Xi_N|\ge\den(\Xi_N)^{-1}$ requires
\begin{equation}
 d_N\log N>A_NN+\log C_N.
 \label{eq:jet-threshold}
\end{equation}
In particular, fixed $A_N=A>0$, fixed $d_N=d$, and subexponential $C_N$ can never suffice.
If $A_N$ stays bounded away from zero and $C_N=e^{o(N)}$, then necessarily
\begin{equation}
 d_N\ge (A_N+o(1))\frac{N}{\log N}.
 \label{eq:jet-order}
\end{equation}
\end{proposition}

\begin{proof}
A contradiction requires the analytic upper bound to be smaller than the arithmetic lower
bound:
\[
 C_NN^{-d_N}<e^{-A_NN}.
\]
Taking logarithms gives~\eqref{eq:jet-threshold}; the remaining statements are immediate.
\end{proof}

\begin{warning}
For actual variable-degree determinants, $A_N$ usually grows with the degree because every
additional node carries an exponential height cost.  Thus~\eqref{eq:jet-order} is only a
necessary benchmark, not a promise that order $N/\log N$ is enough.  The supplied report's
large determinant audits analyze precisely this extra coefficient and denominator growth.
\end{warning}

\subsection{Why the 2026 one-power win disappears}

On the rational grid $n/q$, a zero of order $2k+1$ contributes $q^{-(2k+1)}$, while the
cleared denominator in Theorem~\ref{thm:LMS} is only $q^{2k}$.  The exponent gap of one locks
the residual to zero.

On the conditional orbit, a return of index $N$ is only polynomially close to a previous
point, whereas the denominator is $\asymp W^N$ or $\asymp Z^N$.  Any fixed-order Taylor gain
is therefore polynomial, while the rational lower bound is exponentially small.  The
inequality points in the wrong direction by an exponential margin.  This is the exact reason
the new rational-map theorem does not cross the Alaoglu--Erd\H{o}s wall.

\section{Further proof attempts pushed to explicit failure points}
\label{sec:attempts}

This section records several routes explored after the supplied report.  The purpose is not
to accumulate reformulations, but to identify the first unjustified inference in each route.

\subsection{Route A: use a smooth rational parametrization}

Because the orbit is countable and dense, it has the same order type as $\Q\cap I$; an order
isomorphism extends to an increasing homeomorphism.  One might hope to choose the
homeomorphism smoothly and then invoke rational-map rigidity.

\emph{Exact failure.}  Smoothness is irrelevant to the counting obstruction.
Corollary~\ref{cor:no-tame-conjugacy} forces denominator growth
$\exp(cq^2)$ for every injective rational parametrization.  This is vastly larger than the
polynomial bounds in all known finite-smoothness rigidity theorems.

\subsection{Route B: extend rationality from a dense set}

The identity theorem says that two analytic functions agreeing on a dense set agree
everywhere.  It does not say that an analytic function taking rational values on a dense set
is a rational function.  Transcendental entire functions mapping $\Q$ into $\Q$ have been
known since the nineteenth century, and much stronger modern constructions exist.

\emph{Exact failure.}  Rationality is not a closed condition in the real topology.  A limit
of rational values can be arbitrary.  The missing input must control heights or denominators,
not merely density.

\subsection{Route C: rational Pad\'e approximation at an orbit point}

At a rational graph point $(a,b)$, one can choose $R\in\Q(t)$ with $R(a)=b$ and hope for high
contact with $t^\vartheta$.

\emph{Exact failure.}  Proposition~\ref{prop:transverse} gives contact order exactly one.
A higher Pad\'e order requires coefficients involving $\vartheta$; then evaluation at a
rational orbit point is no longer rational, and the arithmetic subunit argument disappears.
Bivariate rational polynomials can attain higher algebraic multiplicity, but
Theorem~\ref{thm:fixed-contact} shows that every fixed order remains exponentially
underpowered.

\subsection{Route D: force $p$-adic continuity}

The hypotheses define a group homomorphism on the rank-two subgroup
\[
 \langle2,3\rangle\subset\Q_{>0},
 \qquad 2^a3^b\longmapsto m^an^b.
\]
If this homomorphism extended continuously to a suitable $p$-adic closure, $p$-adic
logarithms would impose an additional linear relation and might create a mixed
archimedean--nonarchimedean determinant.

\emph{Exact failure.}  A homomorphism from a dense countable subgroup of a compact
$p$-adic group need not be continuous.  Integrality on the positive semigroup only gives
$p$-adic boundedness there; it gives no control on sequences with signed exponents
$2^{a_j}3^{b_j}\to1$.  The required assertion is therefore a new continuity theorem, not a
formal consequence of integrality.

A proof-producing version of this route must establish the implication
\begin{equation}
 2^{a_j}3^{b_j}\to1\text{ in }\Q_p
 \quad\Longrightarrow\quad
 m^{a_j}n^{b_j}\to1\text{ in }\Q_p
 \label{eq:padic-continuity-target}
\end{equation}
for at least one finite place and for all integer sequences $(a_j,b_j)$.  No argument found
in the supplied report, in the inspected literature, or in the present work proves
\eqref{eq:padic-continuity-target}.

\subsection{Route E: equidistribution and shrinking targets}

The rotation $k\alpha\bmod1$ is equidistributed.  Exact integrality, however, asks for
simultaneous hits on grids whose mesh is exponential in the block height.  Equidistribution
controls fixed intervals and, with Diophantine estimates, some polynomially shrinking
intervals.  It does not force a hit in an exponentially shrinking target.  In fact,
Theorem~\ref{thm:exact-orbit} shows that a counterexample would itself generate exactly the
required exceptional orbit, so generic metric heuristics cannot rule it out.

\subsection{Route F: treat the Matrix Coefficient Conjecture as proved}

For the matrix
\[
 M=\begin{pmatrix}\log2&\log3\\\log m&\log n\end{pmatrix},
\]
the Matrix Coefficient Conjecture would turn singularity into a rational row--column
coefficient relation.  In size two this is equivalent to Four Exponentials
\cite{DasguptaKakde2024}.

\emph{Exact failure.}  The relevant matrix-coefficient statement is a conjecture, not a
proved theorem.  Dasgupta and Kakde prove important structural implications and an analogue
of a Masser theorem, but explicitly leave the implication needed to complete their strategy
open.  Citing the paper as an unconditional determinant-nonvanishing theorem would be a
misreading.

\section{Conditional amplifiers that really would prove the conjecture}
\label{sec:criteria}

The following criteria separate useful targets from equivalent restatements.

\begin{criterion}[A third algebraic direction]\label{crit:third-base}
Suppose the hypotheses $2^x,3^x\in\Z$ imply the existence of a positive algebraic number
$c$ such that
\begin{enumerate}[label=(\roman*),leftmargin=2.2em]
\item $c^x$ is algebraic;
\item $\log c$ is linearly independent over $\Q$ from $\log2,\log3$.
\end{enumerate}
Then $x\in\Z$.
\end{criterion}

\begin{proof}
If $x\notin\Q$, then $1,x$ are $\Q$-independent.  The three logarithms are
$\Q$-independent by assumption.  The Six Exponentials Theorem says that among the six
numbers
\[
 2,3,c,2^x,3^x,c^x
\]
at least one is transcendental, contrary to algebraicity.  Hence $x\in\Q$, and
Proposition~\ref{prop:elementary} gives $x\in\Z$.
\end{proof}

The difficulty is not proving the criterion but manufacturing $c$.  Every base generated
multiplicatively from $2,3,m,n$ remains in the same rank-two logarithmic span.

\begin{criterion}[Denominator-tame thickening]\label{crit:thickening}
Assume a counterexample produces, for some interval $J\subset(1,2)$, a set
$\mathcal T\subset\Q\cap J$ such that
\begin{enumerate}[label=(\roman*),leftmargin=2.2em]
\item $t^\vartheta\in\Q$ for every $t\in\mathcal T$;
\item $\#\{t\in\mathcal T:\Hgt(t),\Hgt(t^\vartheta)\le H\}\gg H^\delta$
      for some $\delta>0$;
\item the denominators obey a uniform polynomial transfer law.
\end{enumerate}
Then determinant or rational-point rigidity methods enter a genuinely stronger regime than
the conditional orbit.
\end{criterion}

This is deliberately stated as an amplifier rather than a theorem.  The conditional orbit
has only $\asymp\log H$ points, and Theorem~\ref{thm:farey-orbit} shows that a mere
reparametrization cannot create the required height density.  A successful proof must create
new rational graph points by an arithmetic operation, not relabel the old ones.

\begin{problem}[Prime-supported sparse rigidity]\label{prob:sparse-rigidity}
Let $0<\alpha<1$ be irrational.  Rule out the simultaneous existence of integers
$r\ge1$, $s\ge0$, $W>1$, and $Z>1$ satisfying
\[
 W\text{ odd},\qquad 3\nmid Z,
 \qquad r+\alpha=\log_2W,
 \qquad s+\alpha=\log_3Z.
\]
Equivalently, rule out a dense cyclic sequence on
$\{(2^t,3^t):0\le t<1\}$ whose reduced coordinate denominators are exactly
\[
 2^{\lfloor k\log_2W\rfloor}
 \quad\text{and}\quad
 3^{\lfloor k\log_3Z\rfloor},
\]
with their floor fluctuations synchronized by the same irrational rotation.
\end{problem}

Problem~\ref{prob:sparse-rigidity} deliberately isolates the extra arithmetic structure that
ordinary rational-point counting ignores.  A useful solution would have to exploit the
synchronized pure-prime denominator laws, rather than merely recast the original logarithmic
equality.

\begin{problem}[One finite-place continuity theorem]\label{prob:padic}
Under the counterexample hypotheses, prove~\eqref{eq:padic-continuity-target} for one prime
$p$ at which the logarithms involved are nondegenerate, and combine the resulting $p$-adic
linear dependence with the real determinant~\eqref{eq:logdet}.
\end{problem}

This route is attractive because it would add a genuinely independent local direction.  It
is also dangerous: continuity is precisely the unproved content, so it must not be hidden
inside notation such as $x\log_p2$ for a real transcendental $x$.

\begin{problem}[Orbit-scale rational rigidity]\label{prob:orbit-rigidity}
Find a rational-map rigidity theorem whose hypotheses use the exact cyclic denominator law
\eqref{eq:Rexact}--\eqref{eq:Sexact} rather than requiring all rational arguments, and whose
conclusion excludes the transcendental power function.  The theorem must remain effective
when the rational-point count is only $c\log H$.
\end{problem}

Theorem~\ref{thm:farey-orbit} shows why a theorem proved by first replacing the orbit with all
rationals cannot work with polynomial heights.  Any successful result must be intrinsically
sparse and prime-sensitive.

\section{A sharpened logical map of the remaining wall}
\label{sec:map}

The following table records what each major route would need beyond currently established
results.

\begin{center}
\begin{tabular}{p{0.27\textwidth}p{0.27\textwidth}p{0.36\textwidth}}
\toprule
Route & Established input & Missing proof-producing step\\
\midrule
Four Exponentials & Exact $2\times2$ logarithmic determinant & Four Exponentials itself, or a special integral case\\
Matrix coefficients & Size-two equivalence and partial strategy & Open matrix-coefficient implication\\
Smooth rational maps & 2026 $C^{2k+1}$ rigidity at denominator $q^k$ & All-rational coverage and rational jet; orbit reparametrization costs $e^{cq^2}$\\
Polynomial / Pad\'e & Exact rational contact and denominator clearing & Variable-complexity gain overcoming exponential height tax\\
Rational-point counting & Exact count $1+\lfloor\log H/\log C_0\rfloor$ & Prime-supported theorem effective at logarithmic count\\
$p$-adic logarithms & Local logarithmic rank tools & Continuity of the real-power homomorphism in one finite place\\
Third-base method & Six Exponentials Theorem & Construct one algebraic exponential outside the rank-two span\\
\bottomrule
\end{tabular}
\end{center}

The common pattern is clear.  The hypotheses provide two algebraic directions and one real
transcendental scaling.  Every successful existing theorem wants either a third algebraic
direction, a rational jet, a polynomially dense rational set, or an independent local
continuity statement.  The conditional orbit supplies none of these automatically.

\section{Conclusion}
\label{sec:conclusion}

The work above does not prove the Alaoglu--Erd\H{o}s conjecture.  It does produce a new and
rigorous continuation of the supplied report:

\begin{enumerate}[leftmargin=2.2em]
\item The compact counterexample orbit has exact heights
\[
 \Hgt(R_k)=W^k,\qquad \Hgt(S_k)=Z^k,
\]
not merely exponential upper and lower bounds.

\item Every injection from all rationals in an interval into that orbit has denominator
distortion
\[
 \exp\!\left(\left(\frac{3|I|\log W}{\pi^2}+o(1)\right)Q^2\right)
 \quad\text{and}\quad
 \exp\!\left(\left(\frac{3|I|\log Z}{\pi^2}+o(1)\right)Q^2\right).
\]
This rules out every polynomially tame reparametrization.

\item The July 2026 smooth rational-map rigidity theorem cannot be applied directly: the
orbit is height-sparse, every tame thickening is impossible by the entropy theorem, and the
power curve has a transcendental first derivative at every rational graph point.

\item Fixed algebraic contact gives only polynomial smallness in a recurrence index, while
rational denominator clearing gives an exponential lower bound.  These estimates are
compatible, so a fixed-order subunit argument cannot close the problem.

\item A genuine proof still requires an amplifier: a third algebraic base, a prime-sensitive
rigidity theorem at logarithmic point density, an independently proved $p$-adic continuity
statement, or a new determinant theorem strong enough to establish the needed size-two
matrix coefficient relation.
\end{enumerate}

The central new lesson is quantitative.  The gap between the known 2026 rational-map
rigidity regime and the Alaoglu--Erd\H{o}s orbit is not a missing constant or one more
finite-difference estimate.  It is an entropy gap between polynomially many rational inputs
and logarithmically many rational graph points, forcing $e^{cQ^2}$ denominator distortion.
Any future proof using rational approximation must exploit the synchronized pure-prime
denominators directly rather than trying to replace the orbit by an ordinary rational grid.

\appendix

\section{Farey counting on a fixed interval}
\label{app:farey}

For completeness, we recall the standard estimate used in
Theorem~\ref{thm:farey-orbit}.  For a fixed interval $I=[a,b]\subset(0,1)$ and a fixed
$q$, the number of $p$ such that $p/q\in I$ and $\gcd(p,q)=1$ is
\[
 |I|\varphi(q)+O(\tau(q)),
\]
where $\varphi$ is Euler's totient and $\tau$ the divisor function.  This follows by
M\"obius inversion:
\[
 \sum_{\substack{aq\le p\le bq\\(p,q)=1}}1
 =\sum_{d\mid q}\mu(d)
   \left(\frac{|I|q}{d}+O(1)\right)
 =|I|\varphi(q)+O(\tau(q)).
\]
Summing over $q\le Q$ and using
\[
 \sum_{q\le Q}\varphi(q)=\frac{3}{\pi^2}Q^2+O(Q\log Q),
 \qquad
 \sum_{q\le Q}\tau(q)=O(Q\log Q),
\]
gives
\[
 \#\mathcal F_I(Q)=\frac{3|I|}{\pi^2}Q^2+O_I(Q\log Q).
\]
Only the quadratic main term is needed for the entropy theorem.

\section{A direct denominator-count formulation}
\label{app:den-count}

Let
\[
 \Delta_k=\max\{\den(R_k),\den(S_k)\},
 \qquad C_0=\max\{W,Z\}.
\]
Theorem~\ref{thm:exact-orbit} gives
\[
 \frac{C_0^k}{3}<\Delta_k<C_0^k.
\]
Therefore the denominator counting function satisfies
\begin{equation}
 1+\left\lfloor\frac{\log D}{\log C_0}\right\rfloor
 \le
 \#\{k\ge0:\Delta_k\le D\}
 \le
 1+\left\lfloor\frac{\log(3D)}{\log C_0}\right\rfloor.
 \label{eq:den-count-direct}
\end{equation}
Thus the orbit has an exact logarithmic denominator entropy with leading constant
$1/\log C_0$.  The Farey set has quadratic height entropy with leading constant
$3|I|/\pi^2$.  Equation~\eqref{eq:conjugacy-distortion} is precisely the inequality obtained
by comparing these two entropies.

\section{Literature note as of August 21, 2026}
\label{app:literature}

The literature check for this continuation focused on developments capable of changing the
logical status of the problem.

\begin{itemize}[leftmargin=2em]
\item The classical three-base result follows from Six Exponentials; replacing three bases
      by two remains the difficult case \cite{SenthilKumarEtAl2018}.
\item Dasgupta--Kakde's Matrix Coefficient Conjecture is explicitly conjectural, and its
      $2\times2$ case is equivalent to Four Exponentials
      \cite{DasguptaKakde2024}.
\item Massold's 2025 work develops new transcendence-degree estimates but does not prove the
      Four Exponentials statement needed here \cite{Massold2025}.
\item Lelis--Moreira--Silva's July 2026 theorem supplies the newest directly relevant
      rational-map rigidity result found in the search.  Its all-rational, rational-jet, and
      polynomial-denominator hypotheses are substantially stronger than the data supplied by
      the compact conditional orbit \cite{LelisMoreiraSilva2026}.
\end{itemize}

No inspected paper supplied an unconditional proof of the two-base conjecture.  This
negative literature statement is not used as a mathematical premise; it explains why the
report distinguishes proved auxiliary theorems from the still-open final implication.

\begin{thebibliography}{99}

\bibitem{AlaogluErdos1944}
L. Alaoglu and P. Erd\H{o}s,
\newblock On highly composite and similar numbers,
\newblock \emph{Transactions of the American Mathematical Society} 56 (1944), 448--469.

\bibitem{Baker1975}
A. Baker,
\newblock \emph{Transcendental Number Theory},
\newblock Cambridge University Press, 1975.

\bibitem{BombieriGubler2006}
E. Bombieri and W. Gubler,
\newblock \emph{Heights in Diophantine Geometry},
\newblock Cambridge University Press, 2006.

\bibitem{DasguptaKakde2024}
S. Dasgupta and M. Kakde,
\newblock Ranks of matrices of logarithms of algebraic numbers II: the Matrix Coefficient
Conjecture,
\newblock arXiv:2408.08178, 2024;
\newblock \url{https://arxiv.org/abs/2408.08178}.

\bibitem{LelisMoreiraSilva2026}
J. Lelis, C. G. Moreira, and E. Silva,
\newblock On $C^k$-functions mapping $\mathbb Q$ into itself and Mahler's problem on
Liouville numbers,
\newblock arXiv:2607.24427, 2026;
\newblock \url{https://arxiv.org/abs/2607.24427}.

\bibitem{Massold2025}
H. Massold,
\newblock Transcendence degrees of fields generated by exponentials of products,
\newblock arXiv:2506.01123, 2025;
\newblock \url{https://arxiv.org/abs/2506.01123}.

\bibitem{PilaWilkie2006}
J. Pila and A. J. Wilkie,
\newblock The rational points of a definable set,
\newblock \emph{Duke Mathematical Journal} 133 (2006), 591--616.

\bibitem{SenthilKumarEtAl2018}
K. Senthil Kumar, R. Thangadurai, and V. Kumar,
\newblock On a problem of Alaoglu and Erd\H{o}s,
\newblock \emph{Resonance} 23 (2018), 749--758.

\bibitem{Waldschmidt2000}
M. Waldschmidt,
\newblock \emph{Diophantine Approximation on Linear Algebraic Groups},
\newblock Springer, 2000.

\bibitem{UnifiedReport2026}
\newblock \emph{Alaoglu--Erd\H{o}s Unified Report},
\newblock supplied research dossier,
\newblock file \texttt{alaoglu-erdos-unified-report.tex}, August 2026.

\end{thebibliography}

\end{document}
